\documentclass{../templateWriting/cdwriting}

\newcommand{\authorinfo}{THE FORUM CENTER - NGUYỄN HOÀNG HUY}
\newcommand{\documenttitle}{TÀI LIỆU CHUYÊN ĐỀ WRITING}
\newcommand{\collectiontitle}{Level 3 — W4 (Jan 19-25)}

\begin{document}

\thispagestyle{firstpage}
\makeworkshopheader{\authorinfo}{\documenttitle}{\collectiontitle}
\pagestyle{otherpages}

\workshopprompt{The graphs show figures relating to hours worked and stress levels amongst professionals in eight groups.}

Describe the information shown to a university or collage lecturer.

\cdsection{PHẦN I: PHÂN TÍCH DỮ LIỆU VÀ XÂY DỰNG CHIẾN LƯỢC}

\cdstep{1. Giải mã Biểu đồ \textit{(Data Decoding)}}

Để đạt Band 7.0+, học viên không được nhìn hai biểu đồ này một cách riêng biệt. Phải tìm ra mối liên hệ giữa chúng.
\begin{itemize}
\item \cdkv{Loại biểu đồ}{Kết hợp (Mixed Charts).}
\begin{itemize}
\item \cdkv{Biểu đồ 1}{Cột dọc (Vertical Bar Chart) -> Số giờ làm việc.}
\item \cdkv{Biểu đồ 2}{Tròn (Pie Chart) -> Tỷ lệ bị bệnh do stress.}
\end{itemize}
\item \cdkv{Đối tượng so sánh}{8 nhóm nghề nghiệp (Writers, Doctors, Chefs, Lecturers, Lawyers, Movie Producers, Programmers, Businessmen).}
\item \cdkv{Thì (Tense)}{Hiện tại đơn (Present Simple) vì không có năm tháng cụ thể, đây là sự thật hiển nhiên/số liệu khảo sát chung.}
\item \cdkv{Đơn vị tính}{Giờ/tuần (Hours) và Phần trăm (\%).}
\end{itemize}

\cdstep{2. Phân tích Chi tiết Số liệu \textit{(Data Estimation)}}

Ta cần quét nhanh số liệu để tìm sự bất thường (Anomaly) hoặc tương quan (Correlation).

Nghề nghiệp

Giờ làm (Bar Chart)

Tỷ lệ Stress (Pie Chart)

Nhận xét thông minh (Band 7+ Insight)

Businessmen

Cao nhất (\~{}70h)

Thấp (11\%)

\cdkv{Nghịch lý 1}{Làm nhiều nhất nhưng stress không cao nhất.}

Movie Producers

Cao nhì (\~{}60h)

Cao nhì (18\%)

\cdkv{Tương quan thuận}{Làm nhiều \$\textbackslash\{\}rightarrow\$ Stress nhiều.}

Doctors

Cao ba (\~{}50h+)

Cao ba (15\%)

Tương quan thuận.

Lecturers

Thấp nhất (<30h)

Cao nhất (25\%)

\cdkv{Nghịch lý 2}{Làm ít nhất nhưng stress cao nhất.}

Writers/Chefs

Trung bình/Thấp

Thấp (8-10\%)

Tương quan thuận.

\cdstep{3. Xác định Đặc điểm Chính \textit{(Key Features - Overview)}}

Ở Band 7.0+, Overview không được liệt kê máy móc. Bạn cần chỉ ra mối quan hệ giữa hai biểu đồ.

\cdstep{Bước 1: So sánh độ lớn \textit{(Magnitude)}:}
\begin{itemize}
\item Ai làm nhiều nhất? -> Businessmen.
\begin{itemize}
\item Ai stress nhất? -> Lecturers.
\end{itemize}
\end{itemize}

\cdstep{Bước 2: Tìm mối tương quan \textit{(Correlation)}:}
\begin{itemize}
\item Có phải cứ làm nhiều giờ là stress nhiều không? -> KHÔNG.
\begin{itemize}
\item \cdkv{Evidence}{Businessmen làm việc nhiều nhất nhưng mức độ stress chỉ ở mức trung bình. Ngược lại, Lecturers làm ít giờ nhất nhưng lại chịu stress cao nhất.}
\end{itemize}
\end{itemize}

\cdkv{=> Kết luận cho Overview}{Số giờ làm việc và mức độ stress không phải lúc nào cũng tỷ lệ thuận (direct correlation). Sự chênh lệch đáng chú ý nhất nằm ở nhóm Lecturers và Businessmen.}

\cdstep{4. Chiến lược Nhóm thông tin \textit{(Grouping Strategy)}}

Để bài viết mạch lạc (Coherence \& Cohesion), ta không viết lần lượt từng nghề từ trái sang phải. Ta sẽ nhóm theo logic so sánh.
\begin{itemize}
\item \cdblue{\textbf{Cách tiếp cận tối ưu:}}
\begin{itemize}
\item \cdkv{Thân bài 1 (Body 1)}{Tập trung vào Bar Chart (Giờ làm việc). Phân nhóm những người làm nhiều giờ nhất (Businessmen, Movie Producers) và ít nhất (Lecturers, Chefs).}
\item \cdkv{Thân bài 2 (Body 2)}{Tập trung vào Pie Chart (Stress) nhưng liên kết ngược lại với Body 1. Tại đây, ta làm nổi bật sự mâu thuẫn (Contrast) giữa giờ làm và stress của Lecturers và Businessmen.}
\end{itemize}
\end{itemize}

\cdsection{PHẦN II: HƯỚNG DẪN VIẾT BÀI CHI TIẾT}

\cdstep{1. Mở bài \textit{(Introduction)}}
\begin{itemize}
\item \cdkv{Nhiệm vụ}{Paraphrase đề bài.}
\item \cdkv{Kỹ thuật nâng cao}{Dùng từ vựng chính xác về sức khỏe và nghề nghiệp.}
\item \cdblue{\textbf{Bài mẫu:}}
\end{itemize}

"The bar chart compares the average weekly working hours of eight different professions, while the pie chart illustrates the percentage of people in these occupations suffering from stress-related illnesses."

\cdstep{2. Tổng quan \textit{(Overview)}}
\begin{itemize}
\item \cdkv{Nhiệm vụ}{Nêu bật sự thiếu tương quan giữa giờ làm và stress.}
\item \cdblue{\textbf{Bài mẫu:}}
\end{itemize}

"Overall, it is evident that working hours do not necessarily correlate with stress levels. While Businessmen work the longest hours, it is Lecturers who suffer the most from stress-related ailments, despite having the lightest workload in terms of hours."
\begin{itemize}
\item \cdkv{Phân tích}{Cấu trúc "It is [Subject] who..." (Câu chẻ) dùng để nhấn mạnh. Từ vựng "Ailments" (bệnh tật) và "Workload" (khối lượng công việc) giúp nâng điểm từ vựng.}
\end{itemize}

\cdstep{3. Thân bài 1: Phân tích Giờ làm việc \textit{(The Bar Chart)}}
\begin{itemize}
\item \cdkv{Chiến lược}{Mô tả từ cao xuống thấp, nhóm các đối tượng nổi bật.}
\item \cdblue{\textbf{Bài mẫu:}}
\end{itemize}

"Regarding working hours, Businessmen stand out with the heaviest schedule, averaging approximately 70 hours per week. They are followed by Movie Producers and Doctors, who work about 60 and slightly over 50 hours, respectively. In contrast, the remaining professions engage in significantly fewer working hours. Writers, Lawyers, and Programmers all work between 35 and 45 hours, while Lecturers record the lowest figure on the chart, at roughly 28 hours per week."
\begin{itemize}
\item \cdblue{\textbf{Phân tích:}}
\begin{itemize}
\item \cdkv{Stand out with the heaviest schedule}{Thay vì nói "work the most".}
\begin{itemize}
\item \cdkv{Engage in}{Thay cho "work".}
\item \cdkv{Respectively}{Cấu trúc liệt kê chuẩn xác.}
\end{itemize}
\end{itemize}
\end{itemize}

\cdstep{4. Thân bài 2: Phân tích Stress \& Sự tương phản \textit{(The Pie Chart \& Synthesis)}}
\begin{itemize}
\item \cdkv{Chiến lược}{Đây là đoạn quan trọng nhất để lấy điểm 7.0+. Không chỉ mô tả Pie chart, mà phải so sánh nó với dữ liệu ở Body 1.}
\item \cdblue{\textbf{Bài mẫu:}}
\end{itemize}

"Turning to stress-related illnesses, the data presents a striking discrepancy compared to working hours. Lecturers, despite having the shortest work week, account for the highest proportion of stress sufferers at 25\%. Conversely, Businessmen, who work the longest, represent only 11\% of the total cases. Movie Producers are the only group showing a direct link between long hours and high stress, ranking second in both categories with 18\%. The figures for other professions vary, with Doctors and Writers accounting for 15\% and roughly 8\% of the total, respectively."
\begin{itemize}
\item \cdblue{\textbf{Phân tích:}}
\begin{itemize}
\item \cdkv{Striking discrepancy}{Sự khác biệt đáng kinh ngạc (Cụm từ Band 8).}
\begin{itemize}
\item \cdkv{Direct link}{Mối liên hệ trực tiếp.}
\item \cdkv{Conversely}{Trạng từ nối chỉ sự ngược lại.}
\item \cdkv{Account for / Represent}{Các động từ đa dạng để chỉ số phần trăm.}
\end{itemize}
\end{itemize}
\end{itemize}

\cdsection{PHẦN III: TỪ VỰNG \& CẤU TRÚC NÂNG CAO}

\cdstep{1. Từ vựng về Mối quan hệ số liệu \textit{(Correlation \& Contrast)}}

Từ/Cụm từ

Nghĩa \& Cách dùng

Ví dụ áp dụng

Inverse correlation

Tương quan nghịch (1 cái tăng, 1 cái giảm)

There is an inverse correlation between hours and stress for Lecturers.

Discrepancy

Sự chênh lệch/mâu thuẫn

A discrepancy between workload and illness rates.

Direct link / Correlation

Tương quan thuận (Cùng tăng/giảm)

Only Movie Producers show a direct link between the two factors.

Disproportionate

Không cân xứng

The stress level is disproportionate to the hours worked.

\cdstep{2. Từ vựng về Nghề nghiệp \& Sức khỏe}

Từ cơ bản

Từ nâng cao

Work hours

Workload / Schedule / Working time

Stress sickness

Stress-related ailments / Stress-related illnesses

People who are stressed

Sufferers / Victims / Cases

Jobs

Professions / Occupations / Career paths

\cdstep{3. Cấu trúc ngữ pháp "Ăn điểm" \textit{(Grammar Focus)}}
\begin{itemize}
\item \cdkv{Mệnh đề nhượng bộ (Concession)}{Dùng để nhấn mạnh sự mâu thuẫn.}
\begin{itemize}
\item \cdkv{Cấu trúc}{Despite + Noun/V-ing, Clause.}
\item \cdkv{Ví dụ}{"Despite having the lightest workload, Lecturers suffer the most."}
\end{itemize}
\item \cdblue{\textbf{Mệnh đề quan hệ rút gọn (Reduced Relative Clause):}}
\begin{itemize}
\item \cdkv{Cấu trúc}{Noun + V-ing/V-ed...}
\item \cdkv{Ví dụ}{"Businessmen, working the longest hours, represent only 11\%..." (Thay vì "who work").}
\end{itemize}
\item \cdkv{Câu chẻ (Cleft Sentence)}{Nhấn mạnh đối tượng.}
\begin{itemize}
\item \cdkv{Cấu trúc}{It is [Subject] that/who...}
\item \cdkv{Ví dụ}{"It is Lecturers who suffer the most."}
\end{itemize}
\end{itemize}

\end{document}
